\begin{abstract}
Large language models (LLMs) use a key-value cache (KV-cache) to avoid recomputing historical attention states during inference. Longer contexts and higher request concurrency make cache management a major constraint on efficiency. This survey reviews KV-cache infrastructure, covering the mechanisms that represent, store, access, and reuse cached keys and values. We organize the survey into three categories: representation and compression, storage and access, and reuse and execution management. We examine how these techniques address capacity pressure, access bandwidth bottleneck, and the overhead of constructing or restoring cached state. The review compares architectural changes, compression, memory management, offloading, and reuse in terms of their benefits, costs, and applicability. We emphasize that optimizations can shift bottlenecks and introduce processing, communication, or quality trade-offs, motivating coordinated design across model and system components. Finally, we discuss future requirements for ultra-long contexts and long-lived agentic tasks, including efficient history access and state preservation across tool interactions.
\end{abstract}
